Distributed combined cooling, heating and power (CCHP) links electricity, heat, and cold in one investment decision, so installed capacities determine for years how tightly local supply can track time-varying demand \cite{liu2018combined,wu2016multiobjective}. Multi-objective planning often pairs annualized cost with emissions or primary-energy indicators \cite{wang2019multi,deb2002fast}, yet explicit \emph{source--load matching} is still less common as a named planning objective even though mismatch directly drives grid import, renewable spillage, and under-utilized assets \cite{zhang2020comprehensive}.

Existing matching metrics often follow a two-stage logic: reduce each carrier's net-load trajectory to a scalar and then add those scalars. This loses the fact that electricity, heating, and cooling may be \emph{simultaneously} imbalanced at hour $t$. Source--load coordination in CCHP is therefore a \emph{multi-carrier time-series} problem, not only a sum of carrier-wise volatility summaries.

This paper is a \textbf{concept validation}. We introduce the \textbf{Instantaneous Euclidean Matching Index (IEMI)}: at each time step, net imbalances in electricity, heat, and cold form a three-vector, whose Euclidean length is the instantaneous multi-carrier mismatch; IEMI is the time average of that length. We compare IEMI with (i) economic-only planning and (ii) a standard-deviation baseline (Std) on one documented German community case. The emphasis is methodological clarity and a closed planning-to-operation evaluation chain.

\textbf{Contributions.} (1) A vector-based matching formulation that preserves instantaneous cross-carrier coupling. (2) Evidence on the German case that IEMI widens the available cost--matching design space and improves full-year operational indicators once higher budget levels are accessible. (3) A concise clarification that the $L_2$ norm is chosen for geometric meaning rather than as a claim of uniformly stronger penalization than $L_1$.

The remainder of the paper introduces the system and dispatch model, defines IEMI, presents the planning formulation and evaluation protocol, and then reports results on the German benchmark.
